\documentclass[conference]{IEEEtran}
\usepackage{amsmath,amssymb,booktabs,microtype}
\usepackage{url}

\title{Controller-Semantic Transport: Certifying When Robot Action-Space Conversion Preserves What a Controller Actually Executes}

\author{\IEEEauthorblockN{Anonymous Authors}}

\begin{document}
\maketitle

\begin{abstract}
Robot-learning datasets and policies are routinely converted between absolute
and delta actions, joint and task spaces, and controller-specific
representations. Existing conversion utilities often treat this problem as
vector re-encoding, even though a stateful controller maps an action together
with measured state and controller memory to a sequence of low-level commands.
We introduce \emph{Controller-Semantic Transport} (CST), which treats
action-space conversion as a semantic compilation problem. CST represents one
control period by a stateful intermediate representation containing the full
drive-target trace, endpoint goal, and next controller reference. It classifies
a conversion as exact-trace, goal-only, bounded approximation, or
unrepresentable, and for affine controller regions compiles an analytic
morphism $u_t=P u_s+q$ together with a whole-source-action-box boundedness
certificate. On a public manipulation stack, a real delta-to-absolute
conversion defect motivates the formulation. In controlled validation, CST
constructs exact delta-current to absolute transports, detects an endpoint-
equivalent but trace-inequivalent controller pair with an analytically predicted
maximum substep error of $0.075$, verifies exact-family trace identities over
200 randomized contexts, and validates one compiled whole-box morphism over 500
random actions without a runtime optimizer. We intentionally do not claim plant,
contact, or task-success equivalence: CST certifies controller-command
semantics and refuses stronger claims when the target interface cannot represent
them.
\end{abstract}

\section{Introduction}
Action representation is not a cosmetic implementation choice. Large-scale
robot studies show substantial differences between absolute and delta
representations and between joint- and task-space policies
\cite{feng2026action}. Public robot-learning stacks consequently provide
multiple control modes and dataset conversion utilities
\cite{tao2025maniskill3,mimic2021}. Yet conversion correctness is typically
judged by endpoint values, replay success, or representation-level
round trips.

That criterion is incomplete for stateful controllers. Consider two
joint-position controllers that share the same endpoint target $q^*$.
One linearly interpolates from current joint position to $q^*$ over the
simulation substeps; the other holds $q^*$ immediately. Their endpoint goals
match, but the low-level command sequences do not. A target-relative delta
controller introduces another difference: the accepted command updates an
internal reference that changes the semantics of the next action.

We study the following question:
\begin{quote}
When does a source controller action admit a target-controller action that
preserves what the controller actually commands, rather than merely matching an
endpoint representation?
\end{quote}

Our approach, Controller-Semantic Transport (CST), has three design goals.
First, the relevant semantics must include controller memory and substep
commands. Second, conversion must be \emph{certified}: exact conversion,
endpoint-only agreement, approximation, and non-representability must remain
distinct outcomes. Third, when an exact affine relation exists, it should be
compiled once rather than re-solved for every action.

Our current contributions are:
\begin{itemize}
\item a stateful controller IR that exposes drive-target trace, endpoint goal,
and next controller reference;
\item a bounded transport compiler that distinguishes exact-trace,
goal-only, bounded-approximate, and unrepresentable conversions;
\item an analytic whole-action-box morphism certificate for affine controller
regions;
\item a public-stack case study motivated by a real delta-to-absolute
conversion failure, while keeping the production bug fix separate from the
research implementation.
\end{itemize}

\section{Controller-Semantic Transport}
\subsection{Stateful semantics}
For native action $u$, measured state $x$, and controller-owned memory $z$,
CST models one control period as
\begin{align}
\tau &= T_u u + T_x x + T_z z + t,\\
g &= G_u u + G_x x + G_z z + g_0,\\
z^+ &= H_u u + H_x x + H_z z + h,
\end{align}
where $\tau$ stacks the low-level drive targets issued at all simulation
substeps, $g$ is a canonical physical goal, and $z^+$ is the next controller
reference. The present implementation gives exact affine IRs for absolute
joint-position, current-relative delta, and target-relative delta controllers,
with normalized or physical native actions, finite native bounds, and
interpolation on/off.

This is deliberately a controller-command semantics. We do not infer equality
of physical trajectories under dynamics or contact from equality of $\tau$.

\subsection{Action-level transport}
For a concrete source action, CST solves a bounded target problem that jointly
matches $\tau$, $g$, and $z^+$. The certificate returns one of:
\emph{EXACT\_TRACE}, \emph{GOAL\_ONLY},
\emph{BOUNDED\_APPROXIMATION}, or \emph{UNREPRESENTABLE}.
It additionally reports the full trace residual, maximum per-substep mismatch,
endpoint and memory residuals, target rank/nullity, and native-bound margin.

The distinction between exact-trace and goal-only is essential. If the source
interpolates from $0.4$ to $0.5$ over four substeps while the target holds
$0.5$, both have identical endpoint goal, but the command traces are
$[0.425,0.45,0.475,0.5]$ and $[0.5,0.5,0.5,0.5]$. CST therefore returns
GOAL\_ONLY with maximum substep mismatch $0.075$.

\subsection{Whole-action-box compilation}
Action-by-action optimization is unnecessary when both controller regions are
affine. Stack the three semantic observables:
\begin{equation}
M_su_s+c_s=M_tu_t+c_t.
\end{equation}
CST compiles
\begin{equation}
u_t=P u_s+q,\quad
P=M_t^+M_s,\quad
q=M_t^+(c_s-c_t).
\end{equation}
An exact affine morphism requires
\begin{equation}
M_tP=M_s,\qquad M_tq+c_t=c_s.
\end{equation}

CST then certifies the complete source native action box
$u_s\in[\ell_s,h_s]$. For each target coordinate, exact interval propagation
through $Pu_s+q$ yields $[\ell_t',h_t']$. The morphism is globally authorized
only if this interval is contained in the target native bounds. Thus an
algebraically exact conversion can still be rejected because it requires
unrepresentable target commands.

\section{Evidence}
\subsection{Public-stack motivating defect}
A public manipulation simulator exposes a trajectory conversion from normalized
delta joint-position control to absolute joint-position control. A focused
regression on the existing path shows that an unbatched NumPy trajectory action
is passed into a tensor-only runtime scaling helper before simulator replay.
Our minimal production patch makes the path explicit:
\[
u_{\Delta,\mathrm{native}}
\rightarrow \Delta q_{\mathrm{physical}}
\rightarrow q^*_{\mathrm{physical}}
\rightarrow u_{\mathrm{abs,native}}.
\]
The upstream-oriented patch is intentionally separated from CST research code.
Simulator-level task-success recovery remains pending and is not claimed here.

\subsection{Controlled semantic validation}
Table~\ref{tab:evidence} summarizes the current frozen evidence. All completed
entries are CI-backed. Randomized checks are verification of analytic
identities, not estimates of task success.

\begin{table}[t]
\caption{Current CST evidence.}
\label{tab:evidence}
\centering
\small
\begin{tabular}{p{0.43\linewidth}p{0.42\linewidth}}
\toprule
Test & Result \\
\midrule
Delta-current $\rightarrow$ absolute, matched interpolation &
EXACT\_TRACE \\
Interpolate $\rightarrow$ hold, matched endpoint &
GOAL\_ONLY; max substep error $0.075$ \\
Random exact-family contexts &
200/200 trace identities \\
Compiled whole-box morphism &
500/500 random actions, residual $<10^{-9}$ \\
Algebraically exact but bound-infeasible &
Refused by whole-box certificate \\
Host controller implementation parity &
Supported in official PDJointPosController harness \\
Exact simulator-state command trace &
Pending \\
Second software stack &
Pending \\
\bottomrule
\end{tabular}
\end{table}

The whole-box test is stronger than sampling as a decision procedure: sampling
is used only as a regression check after the affine map and interval certificate
have already been derived analytically.

\section{Related Work and Scope}
Feng et al. provide a large-scale empirical study of action-space design and
show that absolute/delta and joint/task-space choices materially change policy
learning and control behavior \cite{feng2026action}. CST addresses a different
question: given already chosen source and target controller interfaces, can a
converter prove their command semantics equivalent for a specific context or
region?

Robot-learning frameworks such as ManiSkill3 and robomimic expose rich
controller and dataset interfaces \cite{tao2025maniskill3,mimic2021}. CST is
intended as a semantic layer around such interfaces, not a replacement for
their controllers or datasets.

Our current affine IR does not establish equivalence under nonlinear IK branch
changes, saturation outside declared action boxes, plant dynamics, contact,
force/torque control, or task success. These are explicit promotion gates rather
than hidden assumptions.

\section{Discussion}
CST suggests a different contract for robot action converters. Instead of
returning a transformed vector unconditionally, a converter can return a
proof obligation and one of four semantic outcomes. Exact conversions can be
compiled into cheap runtime morphisms; endpoint-only conversions remain useful
but must not be described as exact controller equivalence; approximate
conversions expose their residual; and unsupported conversions fail closed.

The IR now matches the host ManiSkill PDJointPosController implementation in a separate CI job that exercises set_action() and before_simulation_step() across controller modes. The immediate next steps are deliberately external: retain a minimal upstream fix for the motivating defect, obtain simulator/task-level evidence, and instantiate the same semantic compiler in a second robot-learning stack. These tests determine whether CST is a general controller-semantic
abstraction or merely a correct model of one controller family.

\bibliographystyle{IEEEtran}
\bibliography{refs}
\end{document}
